% !PRINTED 563
\weberpage{563}
%<<CONT
number of points with integral coordinates that lie in a region
$G_x^{(t)}$ enlarged in the ratio $t^{\frac1n} : 1$, and this
number is just as great as the number of lattice points that lie
in the region $G_\xi^{(t)}$. We thus have the following theorem:

\gesperrt{If $G_\xi$ is any finite region in the space $R_\xi$
and $G_\xi^{(t)}$ the region derived from $G_\xi$ by
enlargement in the linear ratio $t^{\frac1n} : 1$; if $Z_t$
is the number of lattice points lying in the interior of
$G_\xi^{(t)}$, then the volume determined by $G_\xi$ is}

\weq{15}{V_\xi = \sqrt{D}\ \mathrm{Lim}_{t\,=\,\infty} \frac{Z_t}{t}.}

We now place about every lattice point, as centre, an
$n$-dimensional sphere of radius $\frac12\sqrt{M}$. Since $\sqrt{M}$
is the least distance between any two lattice points, no two of
these spheres can have a region in common; they can at most touch
one another in a point, and none of these spheres contains,
besides its centre, a further lattice point. Such a sphere
has, by (10), the volume

\[
K_n \left(\frac{\sqrt{M}}{2}\right)^{n},
\]

and the total volume of all those spheres whose centre lies in the
interior of $G_\xi^{(t)}$ is

\[
V_1^{(t)} = Z_t\, K_n \left(\frac{\sqrt{M}}{2}\right)^{n}.
\]

But this volume $V_1^{(t)}$ is evidently smaller than $V^{(t)}$ —
at least if we disregard those spheres which project partly out of
the region $G^{(t)}$, and whose total volume, divided by $t$,
converges to zero as $t$ increases indefinitely — because its
region fills only a part of $G^{(t)}$; and consequently $V$ too is
greater than the limiting value of $V_1^{(t)} : t$, thus

\[
V > \mathrm{Lim}\, \frac{Z_t\, K_n}{t} \left(\frac{\sqrt{M}}{2}\right)^{n},
\]

and from this, by (15),

\weq{16}{\sqrt{D} > K_n \left(\frac{\sqrt{M}}{2}\right)^{n}.}

Since it is now easy to determine the volume $K_n$, the problem is
hereby solved, and an upper limit for $M$ found.
